%  LaTeX support: latex@mdpi.com 
%  For support, please attach all files needed for compiling as well as the log file, and specify your operating system, LaTeX version, and LaTeX editor.

%=================================================================
\documentclass[water,article,submit,pdftex,moreauthors]{Definitions/mdpi} 

%=================================================================

% MDPI internal commands
\firstpage{1} 
\makeatletter 
\setcounter{page}{\@firstpage} 
\makeatother
\pubvolume{1}
\issuenum{1}
\articlenumber{0}
\pubyear{2024}
\copyrightyear{2024}
\externaleditor{Academic Editor: Firstname Lastname}
\datereceived{07 March 2024} 
\daterevised{} 
\dateaccepted{} 
\datepublished{} 
\hreflink{https://doi.org/} % If needed use \linebreak
%\doinum{}

%=================================================================
% Add packages and commands here. The following packages are loaded in our class file: fontenc, inputenc, calc, indentfirst, fancyhdr, graphicx, epstopdf, lastpage, ifthen, lineno, float, amsmath, setspace, enumitem, mathpazo, booktabs, titlesec, etoolbox, tabto, xcolor, soul, multirow, microtype, tikz, totcount, changepage, attrib, upgreek, cleveref, amsthm, hyphenat, natbib, hyperref, footmisc, url, geometry, newfloat, caption

\graphicspath{{figures/}} % path to folder with figures
\usepackage{subcaption}
\usepackage{listings}
\usepackage{xcolor}
\usepackage{gensymb} % for \textdegree
\usepackage{tabularx}
\usepackage{amsmath,mathtools,amsthm}
	\setcounter{MaxMatrixCols}{15}
\usepackage{array}
\usepackage{amssymb}
\usepackage{amsfonts}
\usepackage{float}
\usepackage{chngcntr}
\counterwithin*{equation}{section}
\usepackage[super]{nth}
\usepackage{longtable}
\usepackage{tabularx}
\usepackage{multirow}
\usepackage{adjustbox}

%=================================================================
% Full title of the paper (Capitalized)
\Title{Deep Learning Methods for Satellite Image Processing to Monitor Flood Dynamics in Ganges Delta, Bangladesh}

% MDPI internal command: Title for citation in the left column
\TitleCitation{Deep Learning Methods for Satellite Image Processing to Monitor Flood Dynamics in Ganges Delta, Bangladesh}

% Author Orchid ID: enter ID or remove command
\newcommand{\orcidauthorA}{0000-0002-5759-1089} % Polina Add \orcidA{} behind the author's name

% Authors, for the paper (add full first names)
\Author{Polina Lemenkova $^{1}$*\orcidA{}}

%\longauthorlist{yes}

% MDPI internal command: Authors, for metadata in PDF
\AuthorNames{Polina Lemenkova}

% MDPI internal command: Authors, for citation in the left column
\AuthorCitation{Lemenkova, P.}
% If this is a Chicago style journal: Lastname, Firstname, Firstname Lastname, and Firstname Lastname.

% Affiliations / Addresses (Add [1] after \address if there is only one affiliation.)
\address{%
$^{1}$ \quad Department of Geoinformatics, Faculty of Digital and Analytical Sciences, Universität Salzburg, Schillerstraße 30, A-5020 Salzburg, Austria.
}

% Contact information of the corresponding author
\corres{Correspondence: polina.lemenkova@plus.ac.at; Tel.: +43-677-6173-2772 (P.L.)}

% Current address and/or shared authorship
%\firstnote{Current address: Affiliation 3.} 

% Abstract (Do not insert blank lines, i.e. \\) 
\abstract{Mapping spatial data is essential for monitoring flooded areas, prognosis of hazards and prevention of flood risks. Ganges Delta is the world's largest river delta prone to floods which impact social-natural systems through losses of lives, damage to infrastructure and landscapes. Millions of people living in this region are vulnerable to repetitive foods due to exposure, high susceptibility and low resilience. Cumulative effects from monsoon climate, repetitive rainfalls and tropical cyclones, and hydrogeologic setting in Ganges Delta increase probability of floods. The consequences of floods can be mitigated through monitoring using remote sensing (RS) data, which provide maps of the inundated areas. This study, focusing on the southern segment of Ganges Delta, selected four Landsat 8-9 OLI/TIRS satellite images taken in flood (March) and post-flood (November) periods for analysis of flood extent and landscape changes. This research employs Deep Learning (DL) algorithms of GRASS GIS and modules of qualitative and quantitative analysis. The results presented a series of maps visualised based on the processed satellite images for monitoring floods in Ganges Delta, Bangladesh.}

% Keywords
\keyword{machine learning; deep learning; flood; monsoon; cartography; geoinformatics; remote sensing; satellite image; geospatial analysis; GRASS GIS} 

\PACS{91.10.Da; 91.10.Jf; 91.10.Sp; 91.10.Xa; 96.25.Vt; 91.10.Fc; 95.40.+s; 95.75.Qr; 95.75.Rs}
\MSC{86A30; 86-08; 86A99; 86A04}
\JEL{Y91; Q20; Q24; Q23; Q3; Q01; R11; O44; O13; Q5; Q51; Q55; N57; C6; C61}

\begin{document}

%----------- section ----------->
\section{Introduction}

\subsection{Background: problem formulation}
The repetitive floods in Ganges Delta negatively impact social-natural systems causing losses of lives and damage to infrastructure \cite{JERIN2023}, fragmentation and disruption in natural landscapes \cite{MUKHERJEE2021106961}, salinity intrusion during storms and affected crop agriculture \cite{ABEDIN2023315,ALMASUD2020104443}, and affected sustainability of coastal ecosystems \cite{SARKER2024169718,MAINUDDIN2021105740,HASAN2023101028}. Millions of people living in southern Bangladesh and neighbouring regions of India are vulnerable to repetitive foods due to exposure, high susceptibility and low resilience \cite{SINGHA2020106825,RUMPA2023104047}. The examples of social effects from floods include the disruptions to clean drinking water and electricity, broken transport and communication systems, effects regular education and normal working environment and health care support for population \cite{JERIN2023103851,NAHIN2023e14520}. The long-term effects from flood and its consequences which may last for months  seriously affect the social-economic system of Bangladesh. The coastal area of Ganges River Delta supports approximately 56.7 M of rural population \cite{DAS2021101983} where many villages are affected in the flood-prone areas \cite{AZAD2023100498} which requires proactive and intentional monitoring of flood risks to ensure a broader and more diverse implication of the flood mapping \cite{Alietal}. 

\subsection{Current landscape}

This study presents a case of the Ganges River Delta, located in southern Bangladesh, Figure \ref{fig01}. 

\begin{figure}[H]
    \begin{center}
	\hspace{-35pt}\includegraphics[width=14.0cm]{Fig_01.jpg}
    \end{center}
    \caption{Topographic map of Bangladesh with depicted study area showing the placement of the Landsat data over the Ganges River Delta. Software: GMT version 6.4.0. Map source: author.}
    \label{fig01}
\end{figure}

The Ganges Delta is the world's largest river delta with hundreds of millions people living around. Located in the region with monsoon climate and seasonally changing environmental setting, it is subject to regular floods, occasional tropical cyclones, tidal waves and related inundations \cite{Bhardwaj2022}. Complex hydrology of the Ganges Delta increases the consequences of floods through dense network of streams, rivers and tributaries where water remains for a longer period \cite{Islam}. Alluvial clayey sediments dominating in the riverbed of the contribute to the stagnation of water during flood which accumulated in the connected numerous channels, inner lakes, wetlands and flood plains of the Ganges River Delta \cite{Datta,KHAN201927,Lemenkova202029}. The monsoon climate regulates the repeatability of rainfalls and floods in Ganges Delta which normally augment between the months of March and September. 

\subsection{Gap and motivation}
Currently, the cartographic monitoring of flood in Ganges Delta presents efforts on GIS-based mapping which raised the question of need for operative and accurate data processing for flood prediction and management \cite{Sanyal2023,ISLAM2023167,RANA2023}. Nevertheless, operative monitoring of floods require advanced machine learning techniques and programming approach to support real-time flood prognosis systems through the algorithms of artificial intelligence \cite{ZHAO2021126777}. Identifying landscape patterns within coastal and tidal areas, estuary and delta is challenging due to the obscurity and ambiguity of land cover types. Accurate classification of the satellite images for environmental mapping requires advanced methods of scripting languages which aim at robust detecting of the patterns \cite{YANG2020100413,Lemenkova202322,BEVERIDGE2020104843,SADIQ2023111233,Lemenkova202227}. Existing state-of-the-art methods mostly use the traditional GIS for image processing which may lead to misclassification and inaccurate labelling of pixels. 

The uncertainty of flood prediction in Ganges River Delta has an impact on vulnerability to risk and information among the local population. Hence, preventive prediction rises a question of using reliable data and advances methods for operative monitoring of floods and assessment of environmental sustainability \cite{NICHOLLS2016370}. Notably, no prior study has explored the deep learning (DL) methods available in Geographic Resources Analysis Support System (GRASS) Geographic Information System (GIS) through a comprehensive evaluation of multi-spectral imagery for mapping floods in Ganges Delta. As a response to this gap, this study presents reports such approach which is based on using DL methods of image processing by GRASS GIS scripting software. Implementation of this approach is realised through the MLPClassifier algorithms from the embedded Scikit-Learn library of Python. 

\section{Related works}

Evaluating degree and rates of floods and effectiveness of flood prediction can be performed using approaches of image processing to mapping, i.e., cartographic visualization of pre- and post-flood periods detected from analysis of RS data using machine learning (ML) methods. ML emerge as a fundamental technique of image processing and classification and data analysis in a wide variety of applications including case studies on environmental studies \cite{ROCHA2022107415}, Earth science \cite{LI2023103345,Lemenkova202306,XIONG2021145256}, geological and hazard risk assessment \cite{AOURAGH2023100939} and hydrological monitoring systems \cite{RAFIK2023101569} and vegetation analysis \cite{WU2023110723,Lemenkova202310,DOSSANTOS2022106753}. Such generic approach of ML has been noted earlier \cite{9883389,coasts4010008} with mentioned advantages. These include its high sensitivity of ML which enables to detect variations in the satellite images to evaluate environmental processes, and flexibility of image analysis which can be finely adjusted using diverse techniques of ML and algorithms. For example, existing ML methods include such well-known algorithms as Random Forest (RF), Support Vector Machines (SVM), etc.

%----------- subsection ----------->
\subsection{Deep Learning}

As a branch of Machine learning (ML), Deep learning (DL) is increasingly used in Earth science for automatic processing of geospatial data. The problem remains with the technical approach of such methods, especially in Remote Sensing (RS) data processing. DL approaches can be applied by a variety of tools, using programming libraries, embedded plugins and modules in Geographic Information Systems (GIS) software. Amongst geoinformation software, GRASS GIS proposes effective solutions to satellite image processing with the option to utilise diverse choices of models for classification of RS data, Figure \ref{fig02}.

\begin{figure}[H]
    \begin{center}
%    \setlength{\unitlength}{0.5cm}
	\hspace{-35pt}\includegraphics[width=8.0cm]{Fig_02.jpg}
    \end{center}
    \caption{Conceptual structure of the ML and DL algorithms in geospatial data analysis and image processing by GRASS GIS. Graphical software: Inkscape version 1.2. Scheme source: author.}
    \label{fig02}
\end{figure}

Of the existing ML methods, the most widely used approach is the Deep Learning (DL) with many examples of applications for image processing in environmental studies \cite{10441152,AHMED2024167234,HASSAN2023100399}. DL presents advanced technique for evaluating the environmental dynamics during the climate changes assessed from image analysis \cite{10441576}. There are many categories of DL algorithms which have various setting parameters and functionality for measurements of pixels values in the image and classification. Overall, DL algorithms train the computer model using knowledge derived from the input layers and processed in hidden layers which results in accurate modelling as a result of data processing.

The existing algorithms use Neural Networks (NN) to learn semantic representations of features as indicated by pixels on the images hence derived directly from Earth Observation (EO) data. This enables to evaluate landscapes properties and dynamics of the environmental properties \cite{BISWAS2023e21245}, using such methods as segmentation \cite{Lemenkova202320}, feature extraction, and classification \cite{RAISA2024101128}. Further development of DL methods aims at achieving high precision of the image processing as signals obtained from the spectral reflectances of the land cover types detected on the space borne images, as well as the use of complex programming algorithms, for instance, such as in NN \cite{RUDRA2023e16459}. 

When applied to a series of satellite images, ML methods can effectively discover the properties of landscapes through comparative analysis \cite{9324433,Lemenkova202313}. Moreover, ML enables to quantify the patches of landscapes within the raster images through the computer vision algorithms\cite{10281789,jmse11040871,8899013,Lemenkova202021,9333522} and to determine landscape dynamics over a target period for environmental monitoring \cite{9392684,Lemenkova202324}. Finally, other advantages of the ML approach is that it is a resource- and time-effective method based on advanced programming methods which largely involves scripting techniques. Such approach enables to automate data processing and cartographic analysis as mentioned earlier \cite{10282693,Lemenkova202301,10441504}. Hence, using ML enables to indicate flood risk and detect development of crisis situation in flood-prone areas such as delta of Ganges, Bangladesh. Moreover, such methods are especially effective for  detecting correlation between environmental processes and climate change over time which is especially actual for tropical regions that are subject to monsoon. 

%----------- section ----------->
\section{Materials and Methods}
%----------- subsection ----------->

\subsection{Workflow}
The methodological workflow employed in this study consisted of several steps and processes which aimed at analysis of geospatial information and image processing. It also included diverse data obtained from different sources of information. 

While the purpose of Earth observation (EO) data is to provide information on objects and features visible from space, choice of methods serves varied purposes that accommodate specific DL algorithms and techniques of image processing \cite{4736879,5070195}. Among the available methods, this paper employs the available DL solutions of GRASS GIS for satellite image processing for detecting flooded areas in Ganges River Delta, southern Bangladesh. 

%----------- subsection ----------->
\subsection{Theoretical fundamentals of MLP}
The Multilayer perceptron is a class of DL and a conceptual theoretical background to data analysis performed in this study which should be briefly discussed. A multilayer perceptron (MLP) presents a feedforward Artificial Neural Network (ANN), a branch of DL methods \cite{su151813786}. It consists of fully connected neurons with a nonlinear activation function that is organised in at least three layers, Figure \ref{fig03}. 

\begin{figure}[H]
    \begin{center}
	\includegraphics[width=14.0cm]{Fig_03.jpg}
    \end{center}
    \caption{General methodological scheme for the Deep Learning (DL) approach used for satellite image classification and EO data processing. Software: R, library 'grViz'. Diagram source: author.}
    \label{fig03}
\end{figure}

These layers are being able to distinguish data that is not linearly separable and, for time series of satellite images, the variables that are varying in time. This results in different values of pixels as Digital Number (DN) recognised by the model. Such characteristics is caused by the spatio-temporal variability on spectral reflectance of pixels that constitute the image as raster matrix according to the physical properties of landscape and vegetation patches \cite{rs16010184,HUANG20181915,YIN2021102514}. 

\subsection{Data}
 In this study, we used four Landsat 8-9 Operational Land Imager (OLI) and Thermal Infrared Sensor (TIRS) multispectral satellite images for environmental mapping of Ganges Delta, Bangladesh. The images were pre-projected as original data and had the following technical characteristics: Datum and Ellipsoid World Geodetic System 84 (WGS84), stored in Universal Transverse Mercator (UTM) Zone 46 for southern Bangladesh, and Worldwide Reference System (WRS) Path/Row 137/44. The data were collected at Nadir in day time. The Landsat Collection Category T1, Number 2 for all the scenes. Landsat Station Identifier is LGN. Sensor Identifier is Landsat OLI TIRS for all the images, Data Type L2 is OLI TIRS L2SP. The remaining essential characteristics of the images is summarised in Table \ref{tab01}.

\begin{table}[H]
\footnotesize
\caption{Metadata for the Landsat 8-9 OLI/TIRS images obtained from the USGS} \label{tab01}
\begin{adjustwidth}{-\extralength}{0cm}
    \begin{tabularx}{\fulllength}{llll}
    \toprule
       	\textbf{Data Set Attribute} & \textbf{Attribute Value} & \textbf{Attribute Value} & \textbf{Attribute Value} \\ \hline
        \textbf{Landsat Scene Identifier} & LC81370442023066LGN00 & LC81370442022079LGN00 & LC81370442021076LGN00 \\ \hline
        \textbf{Date Acquired} & \textbf{2023/03/07} & \textbf{2022/03/20} & \textbf{2021/03/17} \\ \hline
        \textbf{Roll Angle} & -1 & 0 & 0 \\ \hline
        \textbf{Start Time} & 2023-03-07 04:24:33 & 2022-03-20 04:24:26.058914 & 2021-03-17 04:24:23.120295 \\ \hline
        \textbf{Stop Time} & 2023-03-07 04:25:05 & 2022-03-20 04:24:57.828914 & 2021-03-17 04:24:54.890294 \\ \hline
        \textbf{Land Cloud Cover} & 3.12 & 0.01 & 0.03 \\ \hline
        \textbf{Scene Cloud Cover L1} & 2.97 & 0.01 & 0.03 \\ \hline
        \textbf{Ground Control Points Model} & 732 & 771 & 776 \\ \hline
        \textbf{Ground Control Points Version} & 5 & 5 & 5 \\ \hline
        \textbf{Geometric RMSE Model} & 5683 & 5615 & 5694 \\ \hline
        \textbf{Geometric RMSE Model X} & 3887 & 3857 & 3585 \\ \hline
        \textbf{Geometric RMSE Model Y} & 4146 & 4081 & 4424 \\ \hline
        \textbf{Processing Software Version} & LPGS\_16.2.0 & LPGS\_15.6.0 & LPGS\_15.4.0 \\ \hline
        \textbf{Sun Elevation L0RA} & 51.70012312 & 56.10505202 & 55.19368850 \\ \hline
        \textbf{Sun Azimuth L0RA} & 134.86314148 & 129.90704446 & 131.01649112 \\ \hline
        \textbf{TIRS SSM Model} & FINAL & FINAL & FINAL \\ \hline
        \textbf{Data Type L2} & OLI\_TIRS\_L2SP & OLI\_TIRS\_L2SP & OLI\_TIRS\_L2SP \\ \hline
        \textbf{Satellite} & 8 & 8 & 8 \\ \hline
        \textbf{Scene Center Lat DMS} & "23°06'46.58""N" & "23°06'45.29""N" & "23°06'46.01""N" \\ \hline
        \textbf{Scene Center Long DMS} & "90°23'29.83""E" & "90°23'14.57""E" & "90°24'53.42""E" \\ \hline
        \textbf{Corner Upper Left Lat DMS} & "24°09'00.29""N" & "24°09'00""N" & "24°09'02.38""N" \\ \hline
        \textbf{Corner Upper Left Long DMS} & "89°13'54.73""E" & "89°13'44.11""E" & "89°15'19.58""E" \\ \hline
        \textbf{Corner Upper Right Lat DMS} & "24°11'21.62""N" & "24°11'21.52""N" & "24°11'22.45""N" \\ \hline
        \textbf{Corner Upper Right Long DMS} & "91°30'23.18""E" & "91°30'12.56""E" & "91°31'48.22""E" \\ \hline
        \textbf{Corner Lower Left Lat DMS} & "22°01'28.45""N" & "22°01'28.20""N" & "22°01'30.32""N" \\ \hline
        \textbf{Corner Lower Left Long DMS} & "89°17'26.70""E" & "89°17'16.22""E" & "89°18'50.22""E" \\ \hline
        \textbf{Corner Lower Right Lat DMS} & "22°03'36.04""N" & "22°03'35.93""N" & "22°03'36.76""N" \\ \hline
        \textbf{Corner Lower Right Long DMS} & "91°31'47.39""E" & "91°31'36.95""E" & "91°33'11.12""E" \\ \hline
        \midrule
        \textbf{Data Set Attribute} & \textbf{Attribute Value} & \textbf{Attribute Value} & \textbf{Attribute Value} \\ \hline
        \textbf{Landsat Scene Identifier} & LC91370442023330LGN00 & LC91370442022327LGN01 & LC81370442021332LGN00 \\ \hline
        \textbf{Date Acquired} & \textbf{2023/11/26} & \textbf{2022/11/23} & \textbf{2021/11/28} \\ \hline
        \textbf{Roll Angle} & 0 & 0 & 0 \\ \hline
        \textbf{Start Time} & 2023-11-26 04:24:51 & 2022-11-23 04:25:01 & 2021-11-28 04:24:54.925489 \\ \hline
        \textbf{Stop Time} & 2023-11-26 04:25:23 & 2022-11-23 04:25:32 & 2021-11-28 04:25:26.695489 \\ \hline
        \textbf{Land Cloud Cover} & 0.02 & 0.20 & 0.40 \\ \hline
        \textbf{Scene Cloud Cover L1} & 0.02 & 0.19 & 0.38 \\ \hline
        \textbf{Ground Control Points Model} & 750 & 782 & 768 \\ \hline
        \textbf{Ground Control Points Version} & 5 & 5 & 5 \\ \hline
        \textbf{Geometric RMSE Model} & 6651 & 6490 & 6697 \\ \hline
        \textbf{Geometric RMSE Model X} & 4185 & 4191 & 4329 \\ \hline
        \textbf{Geometric RMSE Model Y} & 5170 & 4955 & 5110 \\ \hline
        \textbf{Processing Software Version} & LPGS\_16.3.1 & LPGS\_16.2.0 & LPGS\_15.5.0 \\ \hline
        \textbf{Sun Elevation L0RA} & 41.83496249 & 42.43660966 & 41.36291892 \\ \hline
        \textbf{Sun Azimuth L0RA} & 154.44654325 & 154.42300606 & 154.52745495 \\ \hline
        \textbf{TIRS SSM Model} & N/A & N/A & FINAL \\ \hline
        \textbf{Data Type L2} & OLI\_TIRS\_L2SP & OLI\_TIRS\_L2SP & OLI\_TIRS\_L2SP \\ \hline
        \textbf{Satellite} & 9 & 9 & 8 \\ \hline
        \textbf{Scene Center Lat DMS} & "23°06'45.65""N" & "23°06'46.48""N" & "23°06'45.22""N" \\ \hline
        \textbf{Scene Center Long DMS} & "90°22'20.75""E" & "90°23'11.94""E" & "90°24'08.21""E" \\ \hline
        \textbf{Corner Upper Left Lat DMS} & "24°08'48.98""N" & "24°08'50.28""N" & "24°09'01.33""N" \\ \hline
        \textbf{Corner Upper Left Long DMS} & "89°12'51.37""E" & "89°13'44.40""E" & "89°14'37.14""E" \\ \hline
        \textbf{Corner Upper Right Lat DMS} & "24°11'11.26""N" & "24°11'11.80""N" & "24°11'22.06""N" \\ \hline
        \textbf{Corner Upper Right Long DMS} & "91°29'19.50""E" & "91°30'12.67""E" & "91°31'05.70""E" \\ \hline
        \textbf{Corner Lower Left Lat DMS} & "22°01'27.01""N" & "22°01'28.20""N" & "22°01'19.67""N" \\ \hline
        \textbf{Corner Lower Left Long DMS} & "89°16'24.02""E" & "89°17'16.22""E" & "89°18'08.71""E" \\ \hline
        \textbf{Corner Lower Right Lat DMS} & "22°03'35.46""N" & "22°03'35.93""N" & "22°03'26.64""N" \\ \hline
        \textbf{Corner Lower Right Long DMS} & "91°30'44.60""E" & "91°31'36.95""E" & "91°32'29.36""E" \\ \hline
        \bottomrule
    \end{tabularx}
    \end{adjustwidth}
\end{table}

The selection of these data is explained by the reputation and reliability of the Landsat imagery widely used for environmental studies. Besides, the Landsat images are available free of charge, have high repeatability of coverage and high quality. Hence, satellite images present a reliable source of information for detecting the extent of floods and visualising inundated areas using time series analysis. Their applications is reported in many existing studies with a focus on environmental monitoring \cite{Sheonty,Islam2016,Lemenkova202229,Ferdous,Lemenkova202321}. The data were collected from the available freely from the United States Geological Survey (USGS) repository EarthExplorer. The original data are visualised in Figure \ref{fig04}. 

\begin{figure}[H]
  \centering
  \begin{tabular}[b]{c}
    \includegraphics[width=2.9cm]{Fig_04a.jpg} \\
    \small (a): March 2014
  \end{tabular}
  \begin{tabular}[b]{c}
    \includegraphics[width=2.9cm]{Fig_04b.jpg} \\
    \small (b): November 2014
  \end{tabular}
    \begin{tabular}[b]{c}
    \includegraphics[width=2.9cm]{Fig_04c.jpg} \\
    \small (c): March 2021
  \end{tabular}
  \begin{tabular}[b]{c}
    \includegraphics[width=2.9cm]{Fig_04d.jpg} \\
    \small (d): November 2021
  \end{tabular}
   \begin{tabular}[b]{c}
    \includegraphics[width=2.9cm]{Fig_04e.jpg} \\
    \small (e): March 2022
  \end{tabular}
  \begin{tabular}[b]{c}
    \includegraphics[width=2.9cm]{Fig_04f.jpg} \\
    \small (f): November 2022
  \end{tabular}
    \begin{tabular}[b]{c}
    \includegraphics[width=2.9cm]{Fig_04g.jpg} \\
    \small (g): March 2023
  \end{tabular}
  \begin{tabular}[b]{c}
    \includegraphics[width=2.9cm]{Fig_04h.jpg} \\
    \small (h:) November 2023
  \end{tabular}
  \caption{Data sources: Landsat 8-9 OLI/TIRS images on Ganges Delta, Bangladesh, collected during the flood period (March) and post-flood period (November) on years 2014, 2021, 2022 and 2023, used for analysis of flood dynamics. Data source: USGS. Compilation source: author.}
  \label{fig04}
\end{figure}

\subsection{Software}
The workflow used in this research employs a chain of modules that use scripting GRASS GIS software \cite{GRASS_GIS_software} for image processing described in earlier works \cite{Lemenkova202323,technologies11020046}. The methodology of time series analysis includes the unsupervised and supervised classification of multispectral satellite images for detecting changes in landscapes around the Ganges Delta in pre- and post-flood periods. First, the data were imported using the GRASS GIS 'r.import' module; the colour composites were generated using module 'r.composite' and visualised using 'd.rast' and 'd.mon' modules. The map in Figure \ref{fig01} is prepared using Generic Mapping Tools (GMT) version 6.4.0 \cite{Wessel} using raster grid of General Bathymetric Chart of the Oceans (GEBCO). 

%----------- subsection ----------->
\subsection{Implementation}
The following is an outline of the essential details regarding the approaches and techniques of the algorithm that compose the GRASS GIS framework using for image analysis. For details, full codes are presented below and explained with essential details. In all the GRASS GIS modules and functions and post-processing of the satellite images using DL approach, it is necessary to apply the module which implement the required functionality for data processing for each image in terms of the identifiers of functions and parameters of image processing. For this purpose, we adopted and modified the approaches of ML modules (libraries) of GRASS GIS which rely on using the Python's Scikit-Learn library for DL techniques.

%----------- subsection ----------->
\subsection{Image processing}
The unsupervised classification is based on clustering using k-means and Maximum Likelihood embedded in the GRASS GIS. The results of this step are used as training data for supervised classification using deep learning approach. This includes the procedures of clustering that generated signature file and report which was performed using k-means algorithm using 'i.cluster' module. The unsupervised classification was done using the 'i.maxlik' module. The results of this step are presented in Figure \ref{fig02} and provided the seed data and training map for the next step of DL modelling. Deep learning (DL) analysis is performed using algorithm of Multilayer Perceptron (MLPClassifier) for automated classification of the series of satellite images. Technically, the DL approach was realised using the 'r.learn.train' module of GRASS GIS. 

\begin{figure}[H]
	\begin{subfigure}[b]{.45\textwidth}
		\centering
		\includegraphics[width=0.98\textwidth]{Fig_05a.jpg}
		\caption{March 2014}
	\end{subfigure}%
	\begin{subfigure}[b]{.45\textwidth}
		\centering
		\includegraphics[width=0.98\textwidth]{Fig_05b.jpg}
		\caption{November 2014}
	\end{subfigure}
		\hfill
		\vspace{1em}
	\begin{subfigure}[b]{.45\textwidth}
		\centering
		\includegraphics[width=0.98\textwidth]{Fig_05c.jpg}
		\caption{November 2014}
	\end{subfigure}
	\begin{subfigure}[b]{.45\textwidth}
		\centering
		\includegraphics[width=0.98\textwidth]{Fig_05d.jpg}
		\caption{November 2014}
	\end{subfigure}
\caption{Examples of different colour composites created using multispectral bands of Landsat 8 OLI/TIRS scene on March and November 2014 by GRASS GIS 'r.composite' module. Here the band numbers correspond to the following Landsat channels: Band 3 -- Green; Band 4 -- Red; Band 5 -- Near Infrared (NIR); Band 7 -- Shortwave Infrared-2 (SWIR-2).}\label{fig05}
\end{figure}

\begin{figure}[H]
%\hspace{50pt}\makebox[0.90\linewidth][r]{%
	\begin{subfigure}[b]{.50\textwidth}
		\centering
		\includegraphics[width=0.98\textwidth]{Fig_06a.jpg}
		\caption{March 2014}
	\end{subfigure}%
	\begin{subfigure}[b]{.50\textwidth}
		\centering
		\includegraphics[width=0.98\textwidth]{Fig_06b.jpg}
		\caption{November 2014}
	\end{subfigure}
%}
\caption{Image processing of the Landsat 8 OLI/TIRS scene on March and November 2014 using unsupervised classification by clustering and Maximal Likelihood method: \textbf{(a)}: March 2014; \textbf{(b)}: November 2014.}\label{fig06}
\end{figure}

\begin{figure}[H]
%\hspace{50pt}\makebox[0.90\linewidth][r]{%
	\begin{subfigure}[b]{.50\textwidth}
		\centering
		\includegraphics[width=0.98\textwidth]{Fig_07a.jpg}
		\caption{March 2014}
	\end{subfigure}%
	\begin{subfigure}[b]{.50\textwidth}
		\centering
		\includegraphics[width=0.98\textwidth]{Fig_07b.jpg}
		\caption{November 2014}
	\end{subfigure}
%}
\caption{Accuracy assessment of image classification of the Landsat 8 OLI/TIRS scenes: \textbf{(a)}: March 2014; \textbf{(b)}: November 2014.}\label{fig07}
\end{figure}

%----------- section ----------->
\pagebreak
\section{Results}

%----------- subsection ----------->
\subsection{K-means classification outcomes}

Flooded areas in Ganges River Delta detected and mapped using MaxLike methods are presented as a series of the resulting maps in Figure \ref{fig08} and Figure \ref{fig09} and deep learning (DL) is presented in Figure \ref{fig10} and Figure \ref{fig11}. The maps are generated using classified images for March showing the flooding event (in Figure \ref{fig08}) and Figure \ref{fig10}, as well as for November (Figure \ref{fig09} and Figure \ref{fig11}) showing the post-flood landscapes, respectively. 

\begin{figure}[H]
\hspace{50pt}\makebox[0.90\linewidth][r]{%
	\begin{subfigure}[b]{.40\textwidth}
		\centering
		\includegraphics[width=0.98\textwidth]{Fig_08a.jpg}
		\caption{2021}
	\end{subfigure}%
	\begin{subfigure}[b]{.40\textwidth}
		\centering
		\includegraphics[width=0.98\textwidth]{Fig_08b.jpg}
		\caption{2022}
	\end{subfigure}%
	\begin{subfigure}[b]{.40\textwidth}
		\centering
		\includegraphics[width=0.98\textwidth]{Fig_08c.jpg}
		\caption{2023}
	\end{subfigure}%
}
\caption{Flooded landscapes in Ganges Delta, Bangladesh mapped using k-means clustering and MaxLike classification of GRASS GIS applied to the Landsat 8-9 OLI/TIRS images: \textbf{(a)}: March 2021; \textbf{(b)}: March 2022; \textbf{(c)}: March 2023.}\label{fig08}
\end{figure}

\begin{figure}[H]
\hspace{50pt}\makebox[0.90\linewidth][r]{%
	\begin{subfigure}[b]{.40\textwidth}
		\centering
			\includegraphics[width=0.98\textwidth]{Fig_09a.jpg}
		\caption{2021}
	\end{subfigure}%
	\begin{subfigure}[b]{.40\textwidth}
		\centering
		\includegraphics[width=0.98\textwidth]{Fig_09b.jpg}
		\caption{2022}
	\end{subfigure}%
	\begin{subfigure}[b]{.40\textwidth}
		\centering
		\includegraphics[width=0.98\textwidth]{Fig_09c.jpg}
		\caption{2023}
	\end{subfigure}%
}
\caption{Post-flood landscapes in Ganges Delta, Bangladesh mapped using k-means clustering and MaxLike classification of GRASS GIS applied to the Landsat 8-9 OLI/TIRS images: \textbf{(a)}: November 2021; \textbf{(b)}: November 2022; \textbf{(c)}: November 2023.}\label{fig09}
\end{figure}

Maps of floods can be used as a broad base of environmental applications and flood hazard assessment. Their application can range from academia and research, industry and development, to governmental and social services for risk prevention. The findings presented in the study revealed a changing pattern of expanded floods in the post-flood period in Ganges Delta, and inequality in distribution of inundated areas covered by water by consequent years (2021, 2022 and 2023). Scripting models of GRASS GIS with open source available specifications catalyse further innovation in risk management and presents solutions for state-of-the-art mapping using programming approaches and machine learning modules. The following 10 land cover types were identified around the Ganges River Delta, Bangladesh: \emph{1); Water; 2) Wetlands and mudflats; 3) Mangrove forests; 4) Sandy areas; 5) Forests; 6) Croplands; 7) Grasslands; 8) Urban settlements; 9) Orchards; 10) Aquaculture.} The information regarding the land cover types is derived from the existing works \cite{Rahman2020,rs13234799,Paszkowski2021} and adopted to the extent of the current study area.

\begin{table}[H]
  \centering
  \begin{tabularx}{\textwidth}{p{0.9cm}p{0.9cm}p{0.9cm}p{0.9cm}p{0.9cm}p{0.9cm}p{0.9cm}p{0.9cm}p{0.9cm}p{0.9cm}p{0.9cm}}
  \toprule
  \multirow{2}{*}{Year} & \multicolumn{10}{c}{Classes of land cover types in March: Ganges-Brahmaputra River Delta} \\
     & 1 & 2 & 3 & 4 & 5 & 6 & 7 & 8 & 9 & 10 \\
    \midrule
    \textbf{2021} & 717 & 252 & 434 & 758 & 1005 & 1035 & 758 & 744 & 927 & 219 \\
    \textbf{2022} &  749 & 236 & 489 & 803 & 1030 & 1100 & 557 & 744 & 800 & 336 \\
     \textbf{2023} &  707 & 343 & 1055 & 381 & 1032 & 997 & 702 & 649 & 796 & 180 \\
\bottomrule
\end{tabularx}
  \caption{Flood period: estimated classes of land cover types for March period for years 2021, 2022 and 2023 in Ganges River Delta}
  \label{tab02}
\end{table}

\begin{table}[H]
  \centering
  \begin{tabularx}{\textwidth}{p{0.9cm}p{0.9cm}p{0.9cm}p{0.9cm}p{0.9cm}p{0.9cm}p{0.9cm}p{0.9cm}p{0.9cm}p{0.9cm}p{0.9cm}}
  \toprule
  \multirow{2}{*}{Year} & \multicolumn{10}{c}{Classes of land cover types in November: Ganges-Brahmaputra River Delta} \\
     & 1 & 2 & 3 & 4 & 5 & 6 & 7 & 8 & 9 & 10 \\
    \midrule
    \textbf{2021} & 689 & 402 & 579 & 565 & 1228 & 1109 & 571 & 849 & 647 & 208 \\
    \textbf{2022} & 700 & 468 & 365 & 587 & 1066 & 1048 & 683 & 942 & 683 & 326 \\
     \textbf{2023} & 715 & 511 & 231 & 599 & 1351 & 1132 & 718 & 707 & 685 & 310 \\
\bottomrule
\end{tabularx}
  \caption{Post-flood period: estimated classes of land cover types for November period for years 2021, 2022 and 2023 in Ganges River Delta}
  \label{tab02}
\end{table}

%----------- subsection ----------->
\subsection{Deep Learning outcomes}

The supervised learning models of GRASS GIS based on the MLPClassifier demonstrated accurate mapping and refined visualisation functionality, Figure \ref{fig10} and Figure \ref{fig11}. While this method was developed to serve a specific goal of image processing, its purpose is applicable as solutions to environmental monitoring with outcomes of detected flooded areas as presented above. Hence, the research outcomes have provided insights into the changing landscapes of the post-flooded period of Ganges Delta. 

\begin{figure}[H]
\hspace{50pt}\makebox[0.90\linewidth][r]{%
	\begin{subfigure}[b]{.40\textwidth}
		\centering
			\includegraphics[width=0.98\textwidth]{Fig_10a.jpg}
		\caption{2021}
	\end{subfigure}%
	\begin{subfigure}[b]{.40\textwidth}
		\centering
		\includegraphics[width=0.98\textwidth]{Fig_10b.jpg}
		\caption{2022}
	\end{subfigure}%
	\begin{subfigure}[b]{.40\textwidth}
		\centering
		\includegraphics[width=0.98\textwidth]{Fig_10c.jpg}
		\caption{2023}
	\end{subfigure}%
}
\caption{Flooded landscapes in Ganges Delta, Bangladesh mapped using Deep Learning (DL) classification MLPC of GRASS GIS applied to the Landsat 8-9 OLI/TIRS images: \textbf{(a)}: March 2021; \textbf{(b)}: March 2022; \textbf{(c)}: March 2023.}\label{fig10}
\end{figure}

\begin{figure}[H]
\hspace{50pt}\makebox[0.90\linewidth][r]{%
	\begin{subfigure}[b]{.40\textwidth}
		\centering
			\includegraphics[width=0.98\textwidth]{Fig_11a.jpg}
		\caption{2021}
	\end{subfigure}%
	\begin{subfigure}[b]{.40\textwidth}
		\centering
		\includegraphics[width=0.98\textwidth]{Fig_11b.jpg}
		\caption{2022}
	\end{subfigure}%
	\begin{subfigure}[b]{.40\textwidth}
		\centering
		\includegraphics[width=0.98\textwidth]{Fig_11c.jpg}
		\caption{2023}
	\end{subfigure}%
}
\caption{Post-flood landscapes in Ganges Delta, Bangladesh mapped using Deep Learning (DL) classification MLPC of GRASS GIS applied to the Landsat 8-9 OLI/TIRS images: \textbf{(a)}: November 2021; \textbf{(b)}: November 2022; \textbf{(c)}: November 2023.}\label{fig11}
\end{figure}

Time-series analysis of the Ganges Delta using Landsat 8-9 OLI/TIRS dataset enabled by the USGS was used to plots and graphically display the changes in the flood and post-flood periods for southern Bangladesh for three test periods: 2021, 2022 and 2023. The data for 2014 were used as training polygon for deep learning modelling which requires the classification seed. As demonstrated in this paper, GRASS GIS presents an open-source mapping product that shows high potential for its adoption in environmental monitoring and mapping. 

Class separability matrices of land cover classes identified over Ganges River Delta which are presented in the matrices below: Matrix \ref{eq1} shows class separability matrices for March and November 2021; Matrix \ref{eq2} shows class separability matrices for March and November 2022,
 and Matrix \ref{eq3} shows class separability matrices for March and November 2023. Class separability matrices evaluated the assignment of the pixels distributed to various land cover classes over the landscapes of Ganges River Delta, hence it serves as a function of the regularisation parameters used for distinguishability of pixels according to their spectral reflectances. In this way, class separability matrices perform an iterative search for cells within the raster matrix of the image considering a range of values of pixels assigned to different land cover types.

\begin{equation}\label{eq1}
\setlength\arraycolsep{3pt}
\footnotesize
\begin{vmatrix}
    & 1 & 2 & 3 & 4 & 5 & 6 & 7 & 8 & 9 & 10 \\
1 & 0 &  &  &  &  &  &  &  &  &  \\
2 & 1.3 & 0 &  &  &  &  &  &  &  &  \\
3 & 2.2 & 1.0 & 0 &  &  &  &  &  &  &  \\
4 & 3.0 & 1.5 & 0.7 & 0 &  &  &  &  &  &  \\
5 & 3.2 & 1.2 & 1.2 & 1.2 & 0 &  &  &  &  &  \\
6 & 3.4 & 1.9 & 1.3 & 0.7 & 1.1 & 0 &  &  &  &  \\  
7 & 3.6 & 2.4 & 1.7 & 1.2 & 1.6 & 0.7 & 0 &  &  &  \\  
8 & 3.7 & 1.7 & 1.8 & 1.7 & 0.7 & 1.4 & 1.7 & 0 &  &  \\
9 & 3.4 & 1.6 & 2.3 & 2.5 & 1.3 & 2.3 & 2.6 & 1.0 & 0 &  \\   
10 & 3.8 & 2.2 & 2.8 & 2.8 & 2.0 & 2.7 & 3.0 & 1.8 & 1.1 & 0 \\      
\end{vmatrix}
 %
 \hspace{10pt}
 \begin{vmatrix}
    & 1 & 2 & 3 & 4 & 5 & 6 & 7 & 8 & 9 & 10 \\
1 & 0 &  &  &  &  &  &  &  &  &  \\
2 & 1.9 & 0 &  &  &  &  &  &  &  &  \\
3 & 2.6 & 1.0 & 0 &  &  &  &  &  &  &  \\
4 & 2.9 & 2.2 & 1.4 & 0 &  &  &  &  &  &  \\
5 & 3.9 & 2.0 & 0.9 & 1.2 & 0 &  &  &  &  &  \\
6 & 4.6 & 2.4 & 1.4 & 1.5 & 0.6 & 0 &  &  &  &  \\  
7 & 4.2 & 2.8 & 1.9 & 1.3 & 1.3 & 1.0 & 0 &  &  &  \\  
8 & 4.1 & 2.9 & 1.8 & 0.8 & 1.3 & 1.3 & 0.7 & 0 &  &  \\
9 & 3.9 & 3.3 & 2.5 & 1.0 & 2.2 & 2.3 & 1.5 & 1.1 & 0 &  \\   
10 & 3.5 & 3.3 & 2.6 & 1.5 & 2.3 & 2.3 & 1.8 & 1.5 & 0.9 & 0 \\       
 \end{vmatrix} 
\end{equation}

\begin{equation}\label{eq2}
\setlength\arraycolsep{3pt}
\footnotesize
\begin{vmatrix}
   & 1 & 2 & 3 & 4 & 5 & 6 & 7 & 8 & 9 & 10 \\
1 & 0 &  &  &  &  &  &  &  &  &  \\
2 & 1.5 & 0 &  &  &  &  &  &  &  &  \\
3 & 2.4 & 1.0 & 0 &  &  &  &  &  &  &  \\
4 & 3.3 & 1.7 & 0.7 & 0 &  &  &  &  &  &  \\
5 & 3.4 & 1.3 & 1.3 & 1.3 & 0 &  &  &  &  &  \\
6 & 3.8 & 2.2 & 1.3 & 0.7 & 1.2 & 0 &  &  &  &  \\  
7 & 3.6 & 2.4 & 1.6 & 1.1 & 1.5 & 0.6 & 0 &  &  &  \\  
8 & 3.9 & 1.8 & 1.9 & 1.8 & 0.7 & 1.4 & 1.6 & 0 &  &  \\
9 & 3.9 & 1.7 & 2.5 & 2.7 & 1.4 & 2.4 & 2.5 & 1.1 & 0 &  \\   
10 & 3.9 & 2.1 & 2.8 & 2.9 & 2.0 & 2.7 & 2.8 & 1.7 & 1.0 & 0 \\      
\end{vmatrix}
 %
 \hspace{10pt}
 \begin{vmatrix}
    & 1 & 2 & 3 & 4 & 5 & 6 & 7 & 8 & 9 & 10 \\
1 & 0 &  &  &  &  &  &  &  &  &  \\
2 & X & 0 &  &  &  &  &  &  &  &  \\
3 & X & X & 0 &  &  &  &  &  &  &  \\
4 & X & X & X & 0 &  &  &  &  &  &  \\
5 & X & X & X & X & 0 &  &  &  &  &  \\
6 & X & X & X & X & X & 0 &  &  &  &  \\  
7 & X & X & X & X & X & X & 0 &  &  &  \\  
8 & X & X & X & X & X & X & X & 0 &  &  \\
9 & X & X & X & X & X & X & X & X & 0 &  \\   
10 & X & X & X & X & X & X & X & X & X & 0 \\       
 \end{vmatrix} 
\end{equation}

\begin{equation}\label{eq3}
\setlength\arraycolsep{3pt}
\footnotesize
\begin{vmatrix}
    & 1 & 2 & 3 & 4 & 5 & 6 & 7 & 8 & 9 & 10 \\
1 & 0 &  &  &  &  &  &  &  &  &  \\
2 & X & 0 &  &  &  &  &  &  &  &  \\
3 & X & X & 0 &  &  &  &  &  &  &  \\
4 & X & X & X & 0 &  &  &  &  &  &  \\
5 & X & X & X & X & 0 &  &  &  &  &  \\
6 & X & X & X & X & X & 0 &  &  &  &  \\  
7 & X & X & X & X & X & X & 0 &  &  &  \\  
8 & X & X & X & X & X & X & X & 0 &  &  \\
9 & X & X & X & X & X & X & X & X & 0 &  \\   
10 & X & X & X & X & X & X & X & X & X & 0 \\      
\end{vmatrix}
 %
 \hspace{10pt}
 \begin{vmatrix}
    & 1 & 2 & 3 & 4 & 5 & 6 & 7 & 8 & 9 & 10 \\
1 & 0 &  &  &  &  &  &  &  &  &  \\
2 & X & 0 &  &  &  &  &  &  &  &  \\
3 & X & X & 0 &  &  &  &  &  &  &  \\
4 & X & X & X & 0 &  &  &  &  &  &  \\
5 & X & X & X & X & 0 &  &  &  &  &  \\
6 & X & X & X & X & X & 0 &  &  &  &  \\  
7 & X & X & X & X & X & X & 0 &  &  &  \\  
8 & X & X & X & X & X & X & X & 0 &  &  \\
9 & X & X & X & X & X & X & X & X & 0 &  \\   
10 & X & X & X & X & X & X & X & X & X & 0 \\       
 \end{vmatrix} 
\end{equation}
 
%----------- section ----------->
\section{Discussion}

Floods in Ganges River Delta present the most catastrophic problems at the national level in Bangladesh \cite{Szabo}. Preventive mapping of flooded areas presents the societal importance for the social system and maintaining the sustainability of both natural landscapes and social infrastructure. Advanced DL methods for operative mapping of flooded areas result in publicly accessible, modifiable, and distributable open-source cartographic products which aims at prognosis of flood events through machine-based modelling. Hence, operative monitoring and forecasting of natural hazards can effectively reduce the catastrophic consequences of floods in Ganges River Delta. 

Awareness of flood hazards based on maps is an essential tool in addressing environmental challenges through visualization of areas at risks, cartographic display of inundated coastal and prognosis of flood in deltaic areas based on time series analysis. With this regards, using open-source GRASS GIS presents effective retrieval of geo-information using advanced methods of deep learning for accurate satellite image processing. Thus, the displayed time series of maps on Ganges Delta showing floods in March and post-flooded landscapes in November highlighted the benefits achieved by remote sensing data visualization for analysis of flood extent and comparison of climatic effects by recent consecutive years: 2021, 2022 and 2023. 

%----------- section ----------->
\section{Conclusions}

Deep learning in GIS domain enables to increase the accuracy and automation of mapping through the algorithms of artificial intelligence employed in the cartographic workflow. Hence, this research has demonstrated an example of using such approach in cartography and remote sensing data processing. Thus, it has led to a greater understanding of the importance of using advanced methods of machine learning for environmental monitoring and mapping. The application of the ANN, presented here as a case of Multilayer perceptron algorithms demonstrates that this technique can be successfully applied in similar studies which require remote sensing data for spatio-temporal analysis.

This study demonstrates the capacity of sequential analysis of the remote sensing data to identify changes in the coastal and wetland landscapes of Ganges River Delta over the last three years (2021-2023) for comparison of flood event with post-flood landscapes, illustrates how the satellite images can help identify extent of floods, and affected areas, and shows what role deep learning techniques of the GRASS GIS can play in the future technology landscapes in similar studies on vulnerability assessment and mapping floods. Importantly, this paper presented a case of flood mapping in Ganges Delta using the integration of open source data and methods: the dataset was obtained from the publicly available repository of the USGS and processed using freely available GRASS GIS software. Such approach is essential for the environmentalists from the developing countries such as Bangladesh, because open source data and tools can motivate researchers to continue and expand research using the proposed techniques. 

This works can be extended for further environmental monitoring of the Ganges Delta by producing maps on other periods using other EO datasets . 

%----------- section ----------->

\institutionalreview{Not applicable.}

\informedconsent{Not applicable.}

\dataavailability{Not applicable} 

\acknowledgments{The authors thank the reviewers for reading and review of this manuscript.}

\conflictsofinterest{The author declares that they have no known competing financial interests neither personal relationships that could have influenced the work reported in this article.} 

\abbreviations{Abbreviations}{
The following abbreviations are used in this manuscript:\\

\noindent 
\begin{tabular}{@{}ll}
ANN & Artificial Neural Networks \\
DL & Deep Learning \\
DN & Digital Number \\
EO & Earth Observation \\
GEBCO & General Bathymetric Chart of the Oceans \\
GMT & Generic Mapping Tools \\
GRASS & Geographic Resources Analysis Support System \\
GIS & Geographic Information System \\
Landsat OLI/TIRS & Landsat Operational Land Imager and Thermal Infrared Sensor \\
ML & Machine Learning \\
MLP & MultiLayer Perceptron \\
NIR & Near Infrared \\ 
RF & Random Forest \\
RS & Remote Sensing \\
SVM & Support Vector Machines \\
SWIR & Shortwave Infrared \\
USGS & United States Geological Survey \\
UTM & Universal Transverse Mercator \\
WGS84 & World Geodetic System 84\\
WRS & Worldwide Reference System  
\end{tabular}
}

%%%%%%%%%%%%%%%%%%%%%%%%%%%%%%%%%%%%%%%%%%
%% Optional
\appendixtitles{no} % Leave argument "no" if all appendix headings stay EMPTY (then no dot is printed after "Appendix A"). If the appendix sections contain a heading then change the argument to "yes".
\appendixstart
\appendix
\section[\appendixname~\thesection]{}
\subsection[\appendixname~\thesubsection]{}
The appendix 

\section[\appendixname~\thesection]{}
All appendix sections must be cited in the main text. In the appendices, Figures, Tables, etc. should be labeled, starting with ``A''---e.g., Figure A1, Figure A2, etc.

%%%%%%%%%%%%%%%%%%%%%%%%%%%%%%%%%%%%%%%%%%
\begin{adjustwidth}{-\extralength}{0cm}
%\printendnotes[custom] % Un-comment to print a list of endnotes

\reftitle{References}
%\nocite{*}

%=====================================
% References, variant A: external bibliography
%=====================================
\bibliography{Ganges.bib}

%%%%%%%%%%%%%%%%%%%%%%%%%%%%%%%%%%%%%%%%%%
\end{adjustwidth}
\end{document}
